\documentclass[12pt]{article}
\usepackage[utf8]{inputenc}
\usepackage[T1]{fontenc}
\usepackage[margin=1in]{geometry}
\usepackage{setspace}
\usepackage{xcolor}
\usepackage[colorlinks=true, urlcolor=blue!60!black, linkcolor=blue!60!black]{hyperref}
\setlength{\parindent}{0em}
\setlength{\parskip}{0.6em}
\renewcommand{\baselinestretch}{1.0}
\begin{document}
\begin{center}
{\LARGE \textit{District of Columbia v. Heller}, 554 U.S. 570 (2008)\par}
\vspace{0.5em}
{\large Case Brief --- Two-Page Edition\par}
\end{center}

\section{Memory Jogger}
The Second Amendment protects an individual right to keep and bear arms, unconnected with militia service; D.C.'s handgun ban and trigger-lock mandate fell (\href{https://supreme.justia.com/cases/federal/us/554/570/}{554 U.S.\ 570}).

\section{Detailed Case Facts}
The District of Columbia effectively banned civilian handguns by refusing to register them, barred carrying an unlicensed handgun even at home, and required other firearms to be kept ``unloaded and disassembled or bound by a trigger lock or similar device.'' Dick Heller, a special police officer, was denied registration of a handgun for home defense.

\section{Procedural History}
The district court dismissed; the D.C.\ Circuit found Heller alone had standing and reversed, 2--1. The Supreme Court granted certiorari.

\section{Judicial Votes}
Affirmed, 5--4. Scalia wrote for Roberts, Kennedy, Thomas, and Alito. Stevens and Breyer each dissented, each joined by the other dissenters (Souter and Ginsburg).

\section{Holding}
The Amendment guarantees a pre-existing individual right whose core is self-defense in the home. A ban on handguns, an entire class of arms in common use for self-defense, and a requirement that home firearms be inoperable fail under any standard of scrutiny. The Court flagged ``presumptively lawful'' regulations, such as felon bans, as untouched.

\section{Analysis of Opinions}
Scalia's originalist reading held that the prefatory militia clause announces a purpose but cannot limit the operative clause; ``the right of the people'' is individual, and ``bear arms'' meant carrying weapons for confrontation. The right was codified to prevent disarmament of the citizen militia, but self-defense was its ``central component.'' \textit{Miller} protects only arms ``in common use at the time.'' Interest balancing was rejected. Stevens read the Amendment as protecting state militias only and faulted the Court for abandoning reliance on \textit{Miller}. Breyer would uphold the ban under an ``interest-balancing inquiry,'' citing urban gun violence and founding-era urban regulations.

\section{Examples: Future Applications}
\textbf{Same-Side:} Chicago's handgun ban fell once the right was incorporated (\textit{McDonald}). \textbf{Opposite-Side:} disarming people under domestic-violence restraining orders was upheld (\textit{Rahimi}). \textbf{Fence-Sitter:} bans on AR-15-style rifles and large magazines pit ``in common use'' against ``dangerous and unusual weapons.''

\section{Critique}
Critics say the opinion narrows ``the people'' to ``law-abiding, responsible citizens'' without historical proof and uses a circular common-use test; Posner and Wilkinson likened it to \textit{Roe}. Originalists praise its method, though some fault its missing standard of review and Part III exceptions.

\section{Key Quotations}
\begin{quote}
``Although we do not undertake an exhaustive historical analysis today of the full scope of the Second Amendment, nothing in our opinion should be taken to cast doubt on longstanding prohibitions on the possession of firearms by felons and the mentally ill, or laws forbidding the carrying of firearms in sensitive places such as schools and government buildings, or laws imposing conditions and qualifications on the commercial sale of arms.''
\end{quote}
\noindent --- Justice Scalia, Opinion of the Court, 554 U.S.\ at 626--627.

\end{document}
